\documentclass[11pt]{article}
\usepackage[T1]{fontenc}
\usepackage[utf8]{inputenc}
\usepackage{lmodern}
\usepackage{amsmath,amssymb,amsthm}
\usepackage{enumitem}
\usepackage[a4paper,margin=1in]{geometry}
\usepackage{hyperref}
\newtheorem{thm}{Theorem}[section]
\theoremstyle{definition}
\newtheorem{dfn}[thm]{Definition}

\title{Instructions to create automata}
\author{}
%\date{\today}

\begin{document}
\maketitle

\section*{Background definitions}

\noindent\emph{The first goal is to write code to list all admissible ltt structures}
(these will give potential nodes of the automaton). The following definitions build up to defining an admissible ltt structure.

We are interested in a specific kind of colored graph that will have single-edge loops and multi-edges. We provide here this definition and then, for brevity, in the future, unless stated otherwise, we only mean this kind of graph:

\begin{dfn}[graph]
By a \emph{graph} we mean a graph in which each of (1)--(4) holds:
\begin{enumerate}[label=\arabic*)]
  \item There is exactly one red vertex; all other vertices are purple.
  \item There is exactly one red edge; all other edges are either black or purple.
        Edges colored red or purple are called \emph{colored edges}.
  \item Multi-edges are allowed, but for each pair of vertices there is at most one
        colored edge connecting that pair.
  \item Single-edge loops are allowed, but only black edges can be loops.
\end{enumerate}
\end{dfn}

In practice, 

\begin{dfn}[notation]
Define
\[
\mathcal{L}=\{x_1,\bar{x}_1,\ldots,x_5,\bar{x}_5\},\qquad \bar{\bar{x}}_j=x_j.
\]
Assume
\[
x_1<\bar{x}_1<\cdots<x_5<\bar{x}_5
\]
is the increasing order on $\mathcal{L}$.
\end{dfn}

\begin{dfn}[distinguished admissible ltt structure]
A \emph{distinguished admissible ltt structure} is a connected graph $G$ with:
\begin{itemize}
  \item 5 pairs of vertices labeled precisely by the elements of $\mathcal{L}$ (duplicate labels are not allowed);
  \item $x_1$ is red and all remaining vertices are purple;
  \item $x_1$ is connected by a red edge to $x_2$;
  \item $\{\bar{x}_1,x_2,\bar{x}_2,\ldots,x_5,\bar{x}_5\}$ is partitioned into
        three sets of size $3$, and vertices in each partition element are pairwise connected by purple edges;
  \item each pair $\{x_j,\bar{x}_j\}$ is connected by a black edge (this black edge can be viewed as an oriented edge from $\{x_j$ to $\bar{x}_j\}$, and then $\bar{x}_j\}$ would notate this edge with a reverse orientation).
\end{itemize}
\end{dfn}

\begin{dfn}[epi-isomorphic, admissible ltt structure]
Two colored graphs $G_1,G_2$ with vertices labeled by $\mathcal{L}$ are
\emph{epi-isomorphic} via bijection $\phi$ iff there is a bijection $\phi:\mathcal{L}\to\mathcal{L}$ such that:
\begin{enumerate}[label=(\roman*)]
  \item for each $x\in\mathcal{L}$, the vertices labeled $x$ and $\phi(x)$ have
        the same color;
  \item for each distinct pair of elements $x,y\in\mathcal{L}$, the vertices labeled
        $\phi(x),\phi(y)$ are connected by an edge of color $C$ iff the vertices
        labeled $x,y$ are connected by an edge of color $C$ (black here is considered to be a color).
\end{enumerate}
An \emph{admissible ltt structure} is any graph epi-isomorphic to a distinguished
admissible ltt structure.
\end{dfn}



\medskip


\section*{Page 2}

Place some (any) order on the admissible ltt structures.
Let $\mathrm{PLTT}$ denote this ordered list.

\paragraph{Step 1.}
Start with a particular ltt structure $G$.
\begin{itemize}
  \item $\mathrm{GRAPHS}=\{G\}$;
  \item $\mathrm{graphs}_1=\{G\}$;
  \item $\mathrm{boxes}=\{\mathrm{graphs}_1\}$;
  \item layer 1 consists only of $G$, i.e. $\mathrm{GRAPHS}_1=\{G\}$.
\end{itemize}

\noindent\emph{[diagram in source image]}

\medskip
\noindent\textbf{Remark.}
$\mathrm{GRAPHS}$ will ultimately be the list of all nodes in the automaton.
\begin{itemize}
  \item each $\mathrm{graphs}_k$ will be a ``box'' of epi-isomorphic nodes;
  \item $\mathrm{boxes}$ will be the collection of these epi-isomorphism classes;
  \item before closing loops where nodes represent the same epi-isomorphism class of
        ltt structures, the automaton looks schematically like:
        $\mathrm{GRAPHS}_1$ (layer 1), $\mathrm{GRAPHS}_2$ (layer 2),
        $\mathrm{GRAPHS}_3$ (layer 3), etc.
\end{itemize}

\section*{Page 3}

Recursively (with respect to $k$), complete in $\mathrm{PLTT}$ order Steps 2--3
for each $G\in \mathrm{GRAPHS}_k$.

\paragraph{Step 2.}
Find the 4 arrows (terminal ltt structures out of $G$).
Let $k$ denote the layer of $G$, i.e. $G\in\mathrm{GRAPHS}_k$.
\begin{itemize}
  \item Input the starting ltt structure $G$ (an admissible ltt structure):
        vertices labeled by $\mathcal{L}$, with vertex colors and edge colors.
  \item Let $d_r$ be the label of the red vertex.
  \item Let $d_p$ be the label of the purple endpoint of the red edge.
  \item There are exactly two vertices connected to the $d_p$-vertex by purple edges:
        \begin{enumerate}[label=\arabic*)]
          \item let $d_1$ be the label of the first;
          \item let $d_2$ be the label of the second;
        \end{enumerate}
        with order inherited from $\mathcal{L}$.
  \item For each of $d_1,d_2$, there are two potential directed edges out of $G$
        (hence one potential terminal ltt structure for each such edge).
\end{itemize}

For $j\in\{1,2\}$:
\begin{itemize}
  \item[(2a)] edge label: $d_r \to d_j\bar{d}_r$.
        The terminal ltt structure $G_{j;a}$ has:
        \begin{itemize}
          \item red vertex $d_j$ and red edge $\{d_r,\bar{d}_j\}$;
          \item precisely these purple edges:
                \begin{itemize}
                  \item a purple edge $\{d_j,d'\}$ whenever $G$ has a colored edge
                        $\{d_r,d'\}$;
                  \item a purple edge $\{d,d'\}$ whenever $G$ has a purple edge
                        $\{d,d'\}$ and neither $d=d_r$ nor $d'=d_r$.
                \end{itemize}
        \end{itemize}
  \item[(2b)] edge label: $d_j \to d_r\bar{d}_j$.
        The terminal ltt structure $G_{j;b}$ has:
        \begin{itemize}
          \item red vertex $d_j$ and red edge $\{d_j,\bar{d}_r\}$;
          \item precisely these purple edges:
                \begin{itemize}
                  \item a purple edge $\{d_r,d'\}$ whenever $G$ has a colored edge
                        $\{d_j,d'\}$;
                  \item a purple edge $\{d,d'\}$ whenever $G$ has a purple edge
                        $\{d,d'\}$ and neither $d=d_r$ nor $d'=d_r$.
                \end{itemize}
        \end{itemize}
\end{itemize}

Let
\[
\mathrm{terminal}(G):=\{G_{1;a},G_{1;b},G_{2;a},G_{2;b}\}
\]
(repetitions allowed).
\begin{itemize}
  \item Remove any graph from $\mathrm{terminal}(G)$ that is not in $\mathrm{PLTT}$.
  \item Each graph in $\mathrm{terminal}(G)$ is in level $k+1$.
  \item $\mathrm{terminal}(G)$ and level $k+1$ both inherit their order from
        $\mathrm{PLTT}$.
\end{itemize}

\section*{Page 4}

\paragraph{Step 3.}
In $\mathrm{PLTT}$ order, apply the following to each
$G'\in\mathrm{terminal}(G)$.

\begin{itemize}
  \item If $G'$ is epi-isomorphic to some $G''\in\mathrm{GRAPHS}$, but
        $G'\neq G''$:
        \begin{itemize}
          \item $\mathrm{GRAPHS}\leftarrow \mathrm{GRAPHS}\cup\{G'\}$;
          \item $\mathrm{graphs}_j\leftarrow \mathrm{graphs}_j\cup\{G'\}$ where
                $G''\in\mathrm{graphs}_j$;
          \item add a ``permutation arrow'' $G'\to G''$ labeled by $\phi$, where
                $\phi$ is the permutation of $\mathcal{L}$ from the epi-isomorphism.
        \end{itemize}

  \item If there exists $G''\in\mathrm{GRAPHS}$ with $G'=G''$:
        add an arrow $G'\to G''$ with map $G'\mapsto G''$.

  \item For the remaining $G'$:
        \begin{itemize}
          \item $\mathrm{GRAPHS}\leftarrow \mathrm{GRAPHS}\cup\{G'\}$;
          \item $\mathrm{graphs}_{i+1}=\{G'\}$ where $\mathrm{graphs}_i\in\mathrm{boxes}$
                and $\mathrm{graphs}_{i+1}\notin\mathrm{boxes}$;
          \item $\mathrm{GRAPHS}_{k+1}\leftarrow \mathrm{GRAPHS}_{k+1}\cup\{G'\}$.
        \end{itemize}
\end{itemize}

\paragraph{Step 4.}
After Steps 2--3 are completed for each $G\in\mathrm{GRAPHS}_k$, check whether
$\mathrm{GRAPHS}_{k+1}=\varnothing$.

If $\mathrm{GRAPHS}_{k+1}=\varnothing$, terminate.

If $\mathrm{GRAPHS}_{k+1}\neq\varnothing$, continue Steps 2--3 for each
$G\in\mathrm{GRAPHS}_{k+1}$.

\paragraph{Step 5.}
Output the automaton as nodes and edges (including ltt structure/map info).
But what we really want is the strongly connected components.

\end{document}
